\section{Arithmetic specialization}
\label{sec:arithmetic}

The arithmetic exchange is already visible on the cubic \(xyz\).  Its two
ordered icosahedra are Galois conjugate over \(\Q(\sqrt5)\); modulo \(11\)
they become two rational Clebsch parents, and the Galois exchange becomes
the nontrivial spinor class.

\subsection{The golden fibre}

Let \(t^2-t-1=0\), and put
\[
\begin{aligned}
 I_t=\{&
 [0:t:1],[0:t:-1],[1:0:t],[-1:0:t],\\
 & [t:-1:0],[-t:-1:0]\}\subset\mathbf P^2.
\end{aligned}
\]
These six points represent the six opposite vertex pairs of an
icosahedron.

\begin{proposition}[Golden fibre]\label{prop:golden-fibre}
\claimid{HC-2}.
Over a field of characteristic different from \(2\) and \(5\), the points
of \(I_t\) lie on \(xyz=0\), have a common nonzero norm, have normalized
squared inner product \(1/5\), and form a six-arc.  The two configurations
over \(\Q\) are therefore defined over
\[
 \Q[t]/(t^2-t-1)=\Q(\sqrt5)
\]
and are exchanged by its nontrivial automorphism.  In Hitchin's normalized
cover of ordered icosahedra they are the complete fibre over \([xyz]\).
\end{proposition}

\begin{proof}
Every displayed vector has one zero coordinate.  Since \(t^2=t+1\), its
squared norm is \(t^2+1=t+2\).  The inner product of two distinct
representatives is \(\pm t\), while
\[
 (t+2)^2=5(t+1)=5t^2.
\]
This gives the normalized squared inner product \(1/5\).  The twenty
three-point determinants are
\(\pm2t\) or \(\pm2t^2\), so none vanishes under the stated hypotheses.
The final assertion uses Hitchin's classification of the two
icosahedral sets on the orthogonal-plane cubic.
\end{proof}

\subsection{The exchanger}

Set
\[
 R=
 \begin{pmatrix}
 1&0&0\\
 0&0&-1\\
 0&1&0
 \end{pmatrix}.
\]
The relation \(1-t=-1/t\) gives, by substitution,
\[
 R(I_t)=I_{1-t},\qquad
 R^2=\operatorname{diag}(1,-1,-1),\qquad R^4=1.
\]
Moreover \(R(xyz)=-xyz\): it fixes the projective cubic and exchanges its
two ordered icosahedra.

\subsection{Reduction modulo eleven}

The roots of \(t^2-t-1\) in \(\F_{11}\) are \(4\) and \(8\).  For the
working Clebsch arc
\[
\begin{aligned}
 A=\{&
 [1:10:0],[1:9:1],[1:4:7],\\
 & [1:8:5],[0:1:4],[1:1:7]\},
\end{aligned}
\]
define
\[
M_4=\begin{pmatrix}1&4&7\\5&10&3\\1&8&1\end{pmatrix},
\qquad
M_8=\begin{pmatrix}4&5&5\\9&10&7\\4&7&10\end{pmatrix}.
\]
Six projective substitutions give
\[
 M_4I_4=A,\qquad M_8I_8=A,
\]
and \(\det M_4=1,\ \det M_8=9\).  Furthermore,
\[
 M_8^{-1}M_4=3R
\]
in \(\operatorname{GL}_3(\F_{11})\).  Thus forgetting the ordered parent identifies
both configurations with the same unmarked Clebsch arc, while their
relative projectivity is the reduction of the golden exchanger.

\begin{proposition}[Spinor specialization]\label{prop:spinor-specialization}
\claimid{HC-4}.
The reduction of \(R\) represents the nontrivial element of
\[
 T_{11}=\operatorname{PGL}_2(11)/\PSL_2(11).
\]
\end{proposition}

\begin{proof}
For \(Q(x,y,z)=x^2+y^2+z^2\), reflection in a vector \(v\) has spinor
class \(Q(v)\).  On the \(yz\)-plane,
\[
 R=s_{e_2}s_{e_2-e_3},
\]
so
\[
 \theta(R)=Q(e_2)Q(e_2-e_3)=2
 \quad\text{in }\F_{11}^{\times}/\F_{11}^{\times2}.
\]
The class of \(2\) is nonsquare modulo \(11\).  The standard
identifications
\[
 \operatorname{SO}_3(11)\cong\operatorname{PGL}_2(11),
 \qquad
 \Omega_3(11)\cong\PSL_2(11)
\]
identify the spinor quotient with \(T_{11}\), proving the claim.
\end{proof}

\paragraph{Integral boundary.}
Over a finite field of characteristic different from \(2\) and \(5\), the
abstract quadratic cover gives
\[
 \#X_f(\F_q)=1+\chi_q(5J(f)).
\]
Its geometric interpretation over finite fields is proved only away
from an unspecified finite set of primes.  This global comparison is
separate from the good reduction at \(11\) of the explicit golden fibre and
matrix \(R\); we do not replace the finite exceptional set by
\(p\ne2,5\).
